% =====================================================================
%  Grundwissen Klasse 12  --  Informatik
%  Aufbau analog zu ab-0-grundwissen.tex (Grundwissen 6 bis 9)
%
%  ZWEI VERSIONEN AUS EINER DATEI:
%    Arbeitsblatt : einfach uebersetzen
%    Loesung      : \loesungtrue setzen (Zeile unten)  ODER
%                   pdflatex "\def\loesungan{}\input{ab-0-grundwissen-12}"
% =====================================================================
\documentclass[11pt,a4paper]{article}
\input{ab-vorlage}
\renewcommand{\themenbereich}{Grundwissen 6 bis 11}
% schmalere Raender, damit Blatt und Loesung auf 4 Seiten passen
\newgeometry{left=14mm,right=14mm,top=21mm,bottom=12mm,headheight=11mm,headsep=3mm,footskip=5mm}
\setlength{\headwidth}{\textwidth}

% ------------------------- Loesungsschalter --------------------------
\newif\ifloesung
\loesungfalse                      % <-- \loesungtrue  = Loesungsblatt
\ifdefined\loesungan\loesungtrue\fi

\makeatletter
\@ifpackageloaded{xcolor}{}{\usepackage{xcolor}}
\makeatother
\usepackage{comment}
\definecolor{lsgrot}{RGB}{190,30,45}

% Bloecke, die nur in einer der beiden Versionen erscheinen
% (verbatim-sicher, daher auch fuer die java-Umgebung geeignet)
\ifloesung
  \includecomment{nurlsg}\excludecomment{nurab}
\else
  \excludecomment{nurlsg}\includecomment{nurab}
\fi

\newlength{\spaceParagraph}
\setlength{\spaceParagraph}{0.5cm}
% \lsg[Breite]{Loesungstext} -- Luecke bzw. eingetragene Loesung
% (waechst mit, falls die Loesung breiter als die Luecke ist)
\newlength{\lsgbreite}
\newcommand{\lsg}[2][3cm]{%
  \ifloesung
    \settowidth{\lsgbreite}{\small #2}%
    \ifdim\lsgbreite<#1\relax\setlength{\lsgbreite}{#1}\fi
    \underline{\makebox[\lsgbreite][c]{\textcolor{lsgrot}{\small #2}}}%
  \else
    \luecke[#1]%
  \fi}

% Aufgabenbox wie in ab-vorlage, aber nicht umbrechbar: breakable-Boxen
% in Minipages werden sonst am Seitenende gestreckt bzw. zerteilt
\renewenvironment{aufgabe}[2][]{\begin{tcolorbox}[enhanced,colback=white,colframe=black,boxrule=0.7pt,arc=3mm,left=2mm,right=2mm,top=1.5mm,bottom=1.5mm,before skip=2.5mm,after skip=2.5mm]\textbf{Aufgabe #2)}\ifstrempty{#1}{}{ \textit{#1}}\par\vspace{0.6mm}}{\end{tcolorbox}}

% \lsgw[Breite]{Loesungstext} -- wie \lsg, aber umbrechend
% (fuer p-Spalten in Tabellen und lange Antworten)
\newcommand{\lsgw}[2][3cm]{%
  \ifloesung\leavevmode{\small\color{lsgrot}#2}\else\luecke[#1]\fi}

% Alle folgenden Felder sind in Blatt und Loesung gleich hoch,
% damit beide Fassungen identisch umbrechen.
\newlength{\zeilenabstand}
\setlength{\zeilenabstand}{6.5mm}

% \lsgzeilen{Anzahl Linien}{Loesungstext}
\newcommand{\lsgzeilen}[2]{\par\vspace{1mm}\noindent
  \begin{minipage}[t][#1\zeilenabstand][t]{\linewidth}
  \ifloesung{\small\color{lsgrot}#2\par}%
  \else\vspace{0pt}\begin{tikzpicture}\path (0,0);
    \foreach \i in {1,...,#1}{\draw[line width=0.3pt] (0,-\i*\zeilenabstand) -- (\linewidth,-\i*\zeilenabstand);}
  \end{tikzpicture}\fi
  \end{minipage}\par}

% \lsgplatz{Hoehe}{Loesungstext} -- nur in der Loesung sichtbar
\newcommand{\lsgplatz}[2]{\par\noindent
  \begin{minipage}[t][#1][t]{\linewidth}\ifloesung{\small\color{lsgrot}#2\par}\fi\end{minipage}\par}

% Feld fester Hoehe; darin nurab (z.B. \codelinien) und nurlsg (Listing)
\newenvironment{lsgfeld}[1]{\par\noindent\begin{minipage}[t][#1][t]{\linewidth}\vspace{0pt}}{\end{minipage}\par}

% \codelinien{Anzahl} -- gerahmtes Feld mit Schreiblinien fuer Code
\newcommand{\codelinien}[1]{\begin{tikzpicture}
  \draw[line width=0.4pt] (0,0) rectangle (\linewidth-0.4pt,-#1*5mm-2mm);
  \foreach \i in {1,...,#1}{\draw[gray!60,line width=0.3pt] (2mm,-\i*5mm) -- (\linewidth-2mm,-\i*5mm);}
\end{tikzpicture}}

\begin{document}
\abtitel{0}{Grundwissen}

% =====================================================================
\textbf{0.1 Objektorientierte Modellierung -- Klassen und Objekte}\par\nopagebreak

\begin{center}
\begin{tikzpicture}[x=1cm,y=1cm,font=\small]
\node[anchor=north west] (kk) at (0,0) {\klassenkarte{Dreieck}{seitenlaenge\\farbe}{setzeFarbe(neu)}};
\node[anchor=south west] at (0,0.15) {\lsg[5cm]{Klassenkarte}};
\node[anchor=north west,draw,rounded corners=3mm,inner sep=0pt] (ok) at (8.5,0) {\begin{tabular}{l}\rule{0pt}{4mm}\textbf{form1: Dreieck}\\[1mm]\hline\rule{0pt}{4mm}seitenlaenge = 20\\farbe = "blau"\\[1mm]\end{tabular}};
\node[anchor=south west] at (8.5,0.15) {\lsg[5cm]{Objektkarte}};
% Klammern Klassenkarte
\draw[decorate,decoration={brace,amplitude=4pt}] ([xshift=2pt]kk.north east) -- ([xshift=2pt]kk.north east|-{(0,-0.55)}) node[midway,right=6pt] {\lsg[3cm]{Klassenname}};
\draw[decorate,decoration={brace,amplitude=4pt}] ([xshift=2pt]kk.north east|-{(0,-0.65)}) -- ([xshift=2pt]kk.north east|-{(0,-1.55)}) node[midway,right=6pt] {\lsg[3cm]{Attribute}};
\draw[decorate,decoration={brace,amplitude=4pt}] ([xshift=2pt]kk.north east|-{(0,-1.65)}) -- ([xshift=2pt]kk.south east) node[midway,right=6pt] {\lsg[3cm]{Methoden}};
% Klammern Objektkarte
\draw[decorate,decoration={brace,amplitude=4pt}] ([xshift=2pt]ok.north east) -- ([xshift=2pt]ok.north east|-{(0,-0.6)}) node[midway,right=6pt] {\lsg[3cm]{Objektname: Klasse}};
\draw[decorate,decoration={brace,amplitude=4pt}] ([xshift=2pt]ok.north east|-{(0,-0.7)}) -- ([xshift=2pt]ok.south east) node[midway,right=6pt] {\lsg[3cm]{Attributwerte}};
\end{tikzpicture}
\end{center}

\begin{merke}[Objektorientierung]\small
Objektorientierung beschreibt das Lösen von Problemstellungen durch Analyse der beteiligten \textbf{Objekte} und deren Zusammenspiel. Eine \textbf{Klasse} ist der Bauplan, ein \textbf{Objekt} eine konkrete Ausprägung (Instanz) dieser Klasse.
\end{merke}

\vspace{\spaceParagraph/2}
\textbf{0.1.1 Punktnotation für Attribute und Methodenaufrufe}\par\nopagebreak

\begin{minipage}[t]{0.46\linewidth}
\vspace{0pt}
\begin{java}
Objektname.Methode(Parameter);
Objektname.Attribut = Attributwert;
\end{java}
\end{minipage}\hfill
\begin{minipage}[t]{0.50\linewidth}
\vspace{0pt}
\begin{java}
tisch1.breite = 50;
tisch1.verschieben(10,10);
tisch1.loeschen();
karol.schritt();
\end{java}
\end{minipage}

\vspace{2mm}
\textbf{0.1.2 Beziehungen und Klassendiagramme}\par\nopagebreak

Objekte können \textbf{Referenzen} auf andere Objekte enthalten (\textbf{Referenzattribute}). Dabei wird die referenzierte Klasse zum \textbf{Datentyp}:

\begin{minipage}[t]{0.52\linewidth}
\vspace{4pt}
\begin{java}
public class Haus {
   private String strassenname;
   private int grundflaeche;
   private Person eigentuemer;
   ...
}
\end{java}
\end{minipage}\hfill
\begin{minipage}[t]{0.44\linewidth}
\vspace{0pt}
\begin{aufgabe}{1}
\textbf{a)} Zeichne das zugehörige \textbf{Klassendiagramm} (ohne \code{String} und \code{int}).\\
\textbf{b)} Nenne fünf weitere Datentypen.
\end{aufgabe}
\lsgzeilen{4}{b) \code{int}, \code{double}, \code{float}, \code{char}, \code{boolean}, \code{long}, \code{String}\par
a) Klassen \code{Haus} und \code{Person}, verbunden durch eine gerichtete Beziehung
\code{eigentuemer} (Kardinalität 1).}
\end{minipage}

\vspace{\spaceParagraph/2}
\textbf{0.1.3 Kommunikation zwischen Objekten}\par\nopagebreak

Objekte kommunizieren, indem sie 1) öffentliche Attributwerte eines anderen Objekts lesen bzw.\ verändern oder 2) öffentliche Methoden eines anderen Objekts aufrufen. Attribute sollten nicht öffentlich sein (\textbf{Datenkapselung}) -- stattdessen \textbf{Getter}- und \textbf{Setter}-Methoden verwenden.

\vspace{1mm}
\textbf{Zugriffsrechte}\quad\small
\begin{tabular}{|c|l|p{8.4cm}|}\hline
\lsg[0.8cm]{+} & \lsg[2cm]{public} & Öffentlicher Zugriff, jeder kann schreiben und lesen.\\\hline
\lsg[0.8cm]{--} & \lsg[2cm]{private} & \lsgw[8cm]{Zugriff nur innerhalb der eigenen Klasse.}\\\hline
\lsg[0.8cm]{\#} & \lsg[2cm]{protected} & \lsgw[8cm]{Zugriff in der eigenen Klasse und in allen Unterklassen.}\\\hline
\end{tabular}\normalsize

% =====================================================================
\vspace{\spaceParagraph}
\textbf{0.2 Bausteine von Algorithmen -- Struktogramme}\par\nopagebreak

\begin{minipage}[t]{0.47\linewidth}
\textbf{Anweisung}\\[1mm]
\begin{tikzpicture}[x=1cm,y=1cm,font=\ttfamily\small]\draw (0,0) rectangle (4,0.7); \node[anchor=west] at (0.15,0.35) {schritt()};\end{tikzpicture}
\end{minipage}\hfill
\begin{minipage}[t]{0.47\linewidth}
\textbf{Sequenz}\\[1mm]
\begin{tikzpicture}[x=1cm,y=1cm,font=\ttfamily\small]
\draw (0,0) rectangle (4.5,2.1); \draw (0,0.7) -- (4.5,0.7); \draw (0,1.4) -- (4.5,1.4);
\node[anchor=west] at (0.15,1.75) {linksDrehen()}; \node[anchor=west] at (0.15,1.05) {schritt()}; \node[anchor=west] at (0.15,0.35) {aufheben()};
\end{tikzpicture}
\end{minipage}

\vspace{2mm}
\textbf{Die bedingte Anweisung}\par\nopagebreak

\begin{minipage}[t]{0.47\linewidth}
\textit{einseitig}\\[1mm]
\begin{tikzpicture}[x=1cm,y=1cm,font=\small]
\draw (0,0) rectangle (5.5,2.2); \draw (0,2.2) -- (2.75,1.2) -- (5.5,2.2); \draw (0,1.2) -- (5.5,1.2); \draw (2.75,0) -- (2.75,1.2); \draw (2.75,1.2) -- (5.5,0);
\node at (2.75,1.85) {\ttfamily\itshape istZiegel()}; \node at (0.6,1.4) {wahr}; \node at (4.9,1.4) {falsch};
\node at (1.35,0.6) {\ttfamily\small aufheben()};
\end{tikzpicture}\\[1mm]
\textbf{Java:}
\begin{java}
if (istZiegel()) {
   aufheben();
   schritt();
}
\end{java}
\end{minipage}\hfill
\begin{minipage}[t]{0.47\linewidth}
\textit{zweiseitig}\\[1mm]
\begin{tikzpicture}[x=1cm,y=1cm,font=\small]
\draw (0,0) rectangle (5.5,2.2); \draw (0,2.2) -- (2.75,1.2) -- (5.5,2.2); \draw (0,1.2) -- (5.5,1.2); \draw (2.75,0) -- (2.75,1.2);
\node at (2.75,1.85) {\ttfamily\itshape istWand()}; \node at (0.6,1.4) {wahr}; \node at (4.9,1.4) {falsch};
\node at (1.35,0.6) {\ttfamily\small linksDrehen()}; \node at (4.1,0.6) {\ttfamily\small schritt()};
\end{tikzpicture}\\[1mm]
\textbf{Java:}
\begin{lsgfeld}{28mm}
\begin{nurab}
\vspace{1mm}\codelinien{5}
\end{nurab}
\begin{nurlsg}
\begin{java}
if (istWand()) {
   linksDrehen();
} else {
   schritt();
}
\end{java}
\end{nurlsg}
\end{lsgfeld}
\end{minipage}

\textbf{Die Wiederholung}

\begin{minipage}[t]{0.47\linewidth}
\textit{bedingte Wiederholung}\\[1mm]
\begin{tikzpicture}[x=1cm,y=1cm,font=\small]
\draw (0,0) rectangle (5.5,1.9); \draw (0.7,0) -- (0.7,1.3) -- (5.5,1.3);
\node[anchor=west] at (0.15,1.6) {Wiederhole solange {\ttfamily\itshape !istWand()}};
\node[anchor=west] at (0.9,0.85) {\ttfamily\small hinlegen()};
\node[anchor=west] at (0.9,0.4) {\ttfamily\small schritt()};
\end{tikzpicture}\\[1mm]
\textbf{Java:}
\begin{lsgfeld}{23mm}
\begin{nurab}
\vspace{1mm}\codelinien{4}
\end{nurab}
\begin{nurlsg}
\begin{java}
while (!istWand()) {
   hinlegen();
   schritt();
}
\end{java}
\end{nurlsg}
\end{lsgfeld}
\end{minipage}\hfill
\begin{minipage}[t]{0.47\linewidth}
\textit{Wiederholung mit fester Anzahl}\\[1mm]
\begin{tikzpicture}[x=1cm,y=1cm,font=\small]
\draw (0,0) rectangle (5.5,1.9); \draw (0.7,0) -- (0.7,1.3) -- (5.5,1.3);
\node[anchor=west] at (0.15,1.6) {Wiederhole \lsg[2.4cm]{5 mal}};
\node[anchor=west] at (0.9,0.85) {\ttfamily\small hinlegen()};
\node[anchor=west] at (0.9,0.4) {\ttfamily\small schritt()};
\end{tikzpicture}\\[1mm]
\textbf{Java:}
\begin{java}
for (int i = 0; i < 5; i++) {
   hinlegen();
   schritt();
}
\end{java}
\end{minipage}

% =====================================================================
\vspace{\spaceParagraph/2}
\textbf{0.3 Feld (Array)}\par\nopagebreak

\begin{itemize}\itemsep0pt\parskip0pt
\item Zusammenfassung mehrerer Variablen oder Objekte \underline{der gleichen Klasse} in einer zusammenhängenden Folge \underline{fester Länge}.
\item Die einzelnen Feldelemente sind durchnummeriert; diese Nummern nennt man \textbf{Index}. Zugriff über Indizes -- \underline{die Nummerierung beginnt bei 0}.
\end{itemize}

\begin{minipage}[t]{0.42\linewidth}
\vspace{0pt}
\begin{java}
int zahl[];
zahl = new int[1000];
zahl[2] = 10;
\end{java}
\end{minipage}\hfill
\begin{minipage}[t]{0.25\linewidth}
\vspace{0pt}\small
\lsg[2.8cm]{Deklaration}\\[2mm]
\lsg[2.8cm]{Initialisierung}\\[2mm]
\lsg[2.8cm]{Wertzuweisung}
\end{minipage}\hfill
\begin{minipage}[t]{0.29\linewidth}
\vspace{0pt}\small
\underline{Hinweis}: Instanziierung bedeutet die Erzeugung eines neuen Objekts (Aufruf des Konstruktors).
\end{minipage}

% =====================================================================
\vspace{\spaceParagraph}
\textbf{0.4 Das E-V-A-Prinzip}\par\nopagebreak

\begin{minipage}[t]{0.58\linewidth}
\vspace{0pt}
\begin{java}
public void bmi(float gewicht, float groesse)
{
   float bmi = gewicht/(groesse*groesse);
   System.out.println("BMI " + bmi);
}
\end{java}
\end{minipage}\hfill
\begin{minipage}[t]{0.38\linewidth}
\vspace{0pt}
\begin{tikzpicture}[x=1cm,y=1cm,font=\small]
\draw (0,2.3) rectangle (5,2.95); \node at (2.5,2.625) {\lsg[3.2cm]{Eingabe}};
\draw[-latex,thick] (2.5,2.3) -- (2.5,1.95);
\draw (2.5,1.475) ellipse (2.5cm and 0.475cm); \node at (2.5,1.475) {\lsg[3.2cm]{Verarbeitung}};
\draw[-latex,thick] (2.5,1.0) -- (2.5,0.65);
\draw (0,0) rectangle (5,0.65); \node at (2.5,0.325) {\lsg[3.2cm]{Ausgabe}};
\end{tikzpicture}
\end{minipage}

\lsgplatz{5mm}{Eingabe: Parameter \code{gewicht} und \code{groesse} -- Verarbeitung: Berechnung von \code{bmi} --
Ausgabe: \code{System.out.println(...)}}

% =====================================================================
\vspace{\spaceParagraph/2}
\textbf{0.5 Vererbung, Polymorphismus und @Override}\par\nopagebreak

Klassen können voneinander \textbf{erben}. Attribute und Methoden der Oberklasse sind dann auch in Unterklassen verfügbar. \underline{Wichtig}: Eine Klasse kann in Java immer nur von \underline{einer} Klasse erben (keine Mehrfachvererbung). Eine hierarchische Vererbung ist jedoch möglich.

\begin{minipage}[t]{0.46\linewidth}
\vspace{0pt}
\begin{center}
\begin{tikzpicture}[x=1cm,y=1cm,font=\small]
\node[anchor=north west,draw,rounded corners=3mm,inner sep=0pt] at (0,0) {\begin{tabular}{l}\rule{0pt}{4mm}\textbf{person1: Person}\\[1mm]\hline\rule{0pt}{4mm}alter = 17\\name = "Mia"\\[1mm]\end{tabular}};
\node[anchor=north west,draw,rounded corners=3mm,inner sep=0pt] at (4.3,0) {\begin{tabular}{l}\rule{0pt}{4mm}\textbf{person2: Schueler}\\[1mm]\hline\rule{0pt}{4mm}alter = 16\\name = "Tim"\\klasse = "11c"\\[1mm]\end{tabular}};
\end{tikzpicture}
\end{center}
\end{minipage}\hfill
\begin{minipage}[t]{0.50\linewidth}
\vspace{0pt}
\begin{java}
public class Schueler extends Person {
   String klasse;
   public Schueler(int _alter,
         String _name, String _klasse) {
      super(_alter, _name);
      klasse = _klasse;
   }
}
\end{java}
\end{minipage}

\vspace{2mm}
\small
\begin{tabular}{|l|p{12.2cm}|}\hline
\code{super} & Aufruf des Konstruktors (bzw.\ einer Methode) der Oberklasse.\\\hline
\code{extends} & Gibt die Oberklasse an, nach der die Unterklasse erstellt wird.\\\hline
\code{@Override} & \lsgw[11.8cm]{Kennzeichnet eine Methode, die eine geerbte Methode der Oberklasse überschreibt.}\\\hline
Polymorphismus & \lsgw[11.8cm]{Gleicher Methodenaufruf, unterschiedliches Verhalten je nach Objektklasse.}\\\hline
\end{tabular}\normalsize

% =====================================================================
\vspace{\spaceParagraph}
\textbf{0.6 Datenbanken}\par\nopagebreak

\begin{merke}\small
Datenbanken dienen der Speicherung und Verwaltung strukturierter Daten. Die einzelnen Datensätze (\lsg[3cm]{Zeilen}) umfassen dabei mehrere Attributwerte (Spalten $\rightarrow$ \lsg[2.8cm]{Attribute}) eines Objekts und werden in \textbf{Tabellen} gespeichert. Um einen Datensatz eindeutig zu identifizieren, benötigt man einen \lsg[3.6cm]{Primärschlüssel}, welcher aus einem oder mehreren Attributen besteht. Um Beziehungen zwischen Tabellen innerhalb einer Datenbank umzusetzen, werden \lsg[3.6cm]{Fremdschlüssel} benötigt.
\end{merke}

\vspace{\spaceParagraph/2}
\textbf{0.6.1 SQL -- }\lsg[7cm]{Structured Query Language}\par\nopagebreak

\small\begin{tabular}{@{}l l l@{}}
\code{SELECT [DISTINCT]} & \code{<Attribut1>, <Attribut2>, ...} & $\leftarrow$ \lsg[5cm]{Spaltenauswahl}\\
\code{FROM} & \code{<Tabellenname>} & $\leftarrow$ \lsg[5cm]{Quelltabelle(n)}\\
\code{WHERE} & \code{<Bedingung>} & $\leftarrow$ \lsg[5cm]{Zeilenauswahl / Filter}\\
\code{ORDER BY} & \code{<Attribut> [ASC | DESC]} & $\leftarrow$ \lsg[5cm]{Sortierung}\\
\end{tabular}\normalsize

\begin{aufgabe}{2}
Erstelle eine Abfrage auf der Tabelle \code{Schueler}, die alle Nachnamen, Vornamen und Geburtstage der Schüler der Klasse 9a ausgibt, welche im Oktober Geburtstag haben. \textit{Bonus}: Sortiere aufsteigend nach den Geburtstagen.
\end{aufgabe}
\begin{lsgfeld}{27mm}
\begin{nurab}
\codelinien{5}
\end{nurab}
\begin{nurlsg}
\begin{java}[aboveskip=0pt]
SELECT nachname, vorname, geburtstag
FROM Schueler
WHERE klasse = '9a' AND geburtstag LIKE '____-10-__'
ORDER BY geburtstag ASC;
\end{java}
{\small\color{lsgrot}Alternativ (SQLite): \code{strftime('\%m', geburtstag) = '10'}\par}
\end{nurlsg}
\end{lsgfeld}

% =====================================================================
\vspace{\spaceParagraph/2}
\textbf{0.7 Graphen}\par\nopagebreak

\ldots\ bestehen aus \lsg[4cm]{Knoten} und \lsg[4cm]{Kanten}.

\vspace{1mm}
\textbf{a)} \lsg[3cm]{Pfad}: Folge von Knoten, wobei aufeinanderfolgende Knoten mit Kanten verbunden sind und keine Kante doppelt abgelaufen wird.

\textbf{b)} Zyklus: \lsg[11cm]{Pfad, bei dem Start- und Endknoten identisch sind.}\\
\hspace*{5mm}$\rightarrow$ zyklischer oder azyklischer (zyklenfreier) Graph

\textbf{c)} \lsg[3.4cm]{gerichteter} Graph: Kanten können entweder nur in eine (Darstellung als Pfeil) oder in beide Richtungen (Darstellung als Strich) durchlaufen werden.

\textbf{d)} Gewichteter / ungewichteter Graph: \lsg[8.4cm]{Kanten tragen Werte, z.\,B. Entfernungen oder Kosten.}

\textbf{e)} Länge eines Pfades: \lsg[9.6cm]{Anzahl der Kanten bzw. Summe der Kantengewichte.}

% =====================================================================
\vspace{\spaceParagraph}
\textbf{0.8 Client-Server-Modell}\par\nopagebreak

\begin{minipage}[t]{0.40\linewidth}
\vspace{0pt}
\begin{center}
\begin{tikzpicture}[x=1cm,y=1cm,font=\scriptsize,
  geraet/.style={draw,rounded corners=1pt,minimum width=1.05cm,minimum height=0.6cm}]
\node[geraet] (c1) at (2.6,2.6) {Client};
\node[geraet] (c2) at (0.3,1.3) {Client};
\node[geraet] (c3) at (4.9,1.3) {Client};
\node[geraet] (c4) at (2.6,0) {Client};
\node[geraet,fill=black!10] (s) at (2.6,1.3) {\textbf{Server}};
\foreach \c in {c1,c2,c3,c4}{
  \draw[-latex,red!70!black]   (\c) to[bend left=14] (s);
  \draw[-latex,green!45!black] (s)  to[bend left=14] (\c);
}
\node[red!70!black,anchor=west] at (-0.3,0) {Anfrage};
\node[green!45!black,anchor=east] at (5.5,0) {Antwort};
\end{tikzpicture}
\end{center}
\end{minipage}\hfill
\begin{minipage}[t]{0.56\linewidth}
\vspace{0pt}
\begin{aufgabe}{3}
Analysiere die Grafik hinsichtlich ihrer Grapheneigenschaften und recherchiere, welche Protokolle und Adressen zur Kommunikation zwischen den Endpoints nötig sind.
\end{aufgabe}
\lsgzeilen{2}{Gerichteter, ungewichteter, zyklenfreier Graph (Stern\-topologie); der Server ist der zentrale Knoten.}
\end{minipage}

Beispiele für Protokolle: \lsg[11cm]{HTTP/HTTPS, FTP, SMTP, IMAP/POP3, DNS, TCP/IP}

Die Adressierung von Geräten erfolgt über die \lsg[2.4cm]{IP-Adresse} und auf Hardwareseite über die \lsg[2.4cm]{MAC-Adresse}.

% =====================================================================
\vspace{\spaceParagraph/2}
\textbf{0.9 Codierung und Verschlüsselung}\par\nopagebreak

\begin{aufgabe}{4}
\textbf{a)} Erkläre zwei verschiedene Verfahren zur Darstellung der gleichen Information mit einer passenden Grafik und vergleiche dein Vorgehen mit deinem Nachbarn.\\
\textbf{b)} Verschlüssele eine deiner Informationsdarstellungen und tausche sie mit einer Person im Kurs aus (nicht in Armreichweite), ohne aufzustehen.\\
\textbf{c)} Erkläre die Anforderungen an Verschlüsselungsverfahren und gehe dabei auf die zwei grundlegend unterschiedlichen Typen ein.
\end{aufgabe}

\lsgzeilen{5}{a) z.\,B. Dezimal- und Binärdarstellung derselben Zahl, ASCII-Code für Text, RGB-/Hex-Code für Farben:
gleiche Information, unterschiedliche Codierung.\par
c) Anforderungen: Die Sicherheit darf nur vom Schlüssel abhängen, nicht vom Verfahren (Kerckhoffs-Prinzip);
großer Schlüsselraum (kein Ausprobieren); keine erkennbaren Muster/Häufigkeiten im Geheimtext;
effiziente Ver- und Entschlüsselung; sichere Schlüsselübergabe.
Typen: \textbf{symmetrisch} -- ein gemeinsamer Schlüssel zum Ver- und Entschlüsseln (z.\,B. Cäsar, AES),
Problem der Schlüsselübergabe. \textbf{Asymmetrisch} -- öffentlicher Schlüssel zum Verschlüsseln,
privater Schlüssel zum Entschlüsseln (z.\,B. RSA), löst das Schlüsselübergabeproblem.}

% =====================================================================
\vspace{\spaceParagraph/2}
\textbf{0.10 Arten von Künstlicher Intelligenz und Maschinellem Lernen}\par\nopagebreak

\begin{minipage}[t]{0.46\linewidth}
\vspace{0pt}
\begin{center}
\begin{tikzpicture}[font=\tiny]
\draw (0,0) ellipse (2.6cm and 1.6cm); \node at (0,1.3) {Informatik};
\draw[fill=black!8] (0,-0.3) ellipse (2.15cm and 1.3cm); \node at (0,0.7) {Data Science};
\draw[fill=black!16] (0,-0.6) ellipse (1.7cm and 1.0cm); \node at (0,0.1) {Künstliche Intelligenz};
\draw[fill=black!26] (0,-0.9) ellipse (1.25cm and 0.7cm); \node at (0,-0.48) {Maschinelles Lernen};
\draw[fill=black!45] (0,-1.15) ellipse (0.8cm and 0.45cm); \node[white,align=center,inner sep=0pt] at (0,-1.15) {Deep\\Learning};
\end{tikzpicture}
\end{center}
\end{minipage}\hfill
\begin{minipage}[t]{0.50\linewidth}
\vspace{0pt}
\textbf{a)} Grenze \textit{schwache} und \textit{starke} KI voneinander ab.
\lsgzeilen{4}{\textbf{Schwache KI}: auf eine spezifische Aufgabe beschränkt, kein Selbstbewusstsein
(z.\,B. Sprachassistenten, Empfehlungssysteme) -- existiert heute.
\textbf{Starke KI}: könnte jede menschliche Aufgabe bewältigen, mit echtem Verständnis -- bisher hypothetisch.}
\end{minipage}

\vspace{1mm}
\begin{minipage}[t]{0.46\linewidth}
\vspace{0pt}
\textbf{b)} Beschrifte das künstliche Neuron.\\[1mm]
\begin{tikzpicture}[x=1cm,y=1cm,font=\scriptsize]
\node (x1) at (0,1.0) {$x_1$};
\node (x2) at (0,0) {$x_2$};
\node[draw,circle,minimum size=1.1cm] (n) at (2.6,0.5) {$\Sigma\,|\,f$};
\node (y) at (5.0,0.5) {Ausgabe};
\draw[-latex] (x1) -- node[above,pos=0.55] {$w_1$} (n);
\draw[-latex] (x2) -- node[below,pos=0.55] {$w_2$} (n);
\draw[-latex] (n) -- (y);
\draw[-latex] (2.6,-0.45) node[right] {$b$} -- (n);
\end{tikzpicture}
\end{minipage}\hfill
\begin{minipage}[t]{0.50\linewidth}
\vspace{0pt}
\lsgzeilen{4}{$x_1, x_2$ = Eingabewerte, $w_1, w_2$ = Gewichte, $b$ = Schwellenwert/Bias.
Das Neuron bildet die gewichtete Summe $\sum w_i x_i + b$ und wendet darauf eine
Aktivierungsfunktion $f$ an; das Ergebnis ist die Ausgabe. Lernen bedeutet Anpassen der Gewichte.}
\end{minipage}

\vspace{1mm}
\textbf{c)} Grenze kurz die drei Arten Maschinellen Lernens voneinander ab.

\lsgzeilen{4}{\textbf{Überwachtes Lernen (Supervised Learning)}: Training mit gelabelten Beispieldaten (Eingabe + gewünschte Ausgabe),
z.\,B. Klassifikation von Bildern. --
\textbf{Unüberwachtes Lernen (Unsupervised Learning)}: nur Eingabedaten ohne Label; das System erkennt selbst Strukturen
und Muster, z.\,B. Clustering von Kundengruppen. --
\textbf{Verstärkendes Lernen (Reinforcement Learning)}: Lernen durch Belohnung und Bestrafung
in einer Umgebung, z.\,B. Spiel- oder Steuerungsstrategien.}

\end{document}